\documentclass[a4paper,landscape]{article}
\usepackage[T1]{fontenc}
\usepackage[margin=5mm]{geometry}
\usepackage{amsmath,amssymb,multicol,xcolor,helvet}
\renewcommand{\familydefault}{\sfdefault}
\pagestyle{empty}
\setlength{\parindent}{0pt}
\setlength{\parskip}{1.2pt}
\setlength{\columnsep}{4mm}
\setlength{\columnseprule}{.2pt}
\setlength{\abovedisplayskip}{2pt}
\setlength{\belowdisplayskip}{2pt}
\definecolor{early}{HTML}{29577A}
\definecolor{late}{HTML}{166451}
\definecolor{alert}{HTML}{9B3B30}
\newcommand{\head}[2][late]{\par\vspace{3pt}\noindent\colorbox{#1}{\parbox{\dimexpr\linewidth-6pt}{\color{white}\bfseries #2}}\par\vspace{1pt}}
\newcommand{\cue}[1]{\par\textbf{#1}\enspace}
\newcommand{\trap}[1]{\par{\color{alert}\textbf{Watch:} #1}}
\newcommand{\E}{\operatorname{E}}
\newcommand{\V}{\operatorname{Var}}
\newcommand{\Cov}{\operatorname{Cov}}
\newcommand{\SE}{\operatorname{SE}}
\newcommand{\xb}{\bar x}
\newcommand{\Xb}{\bar X}
\newcommand{\db}{\bar d}
\newcommand{\R}{\mathbb R}
\newcommand{\col}{\par\vfill\columnbreak}
\begin{document}
\fontsize{7.4}{8.3}\selectfont
\sffamily\raggedright
\setlength{\abovedisplayskip}{2pt}
\setlength{\belowdisplayskip}{2pt}
\setlength{\abovedisplayshortskip}{1pt}
\setlength{\belowdisplayshortskip}{1pt}
\begin{multicols*}{4}
\head[early]{ST2334 FINALS\hfill 1 / 2}
\textbf{AY26/27 S1 | A4, two sides}\par
Blue: first half (25\%); green: second half (75\%).
\head[early]{1\quad Events, counting, conditioning}
\cue{Count first.} Order? Replacement? Equally likely?
\( {}_nP_r=n!/(n-r)!,\quad {n\choose r}=n!/[r!(n-r)!]\).
Ordered with replacement: \(n^r\); repeated objects: \(n!/\prod n_i!\); circle: \((n-1)!\).
\cue{Restrictions.} Together: block, then permute inside. No adjacency: arrange others, choose gaps, permute. Grid paths: \({a+b\choose a}\); through a point: multiply subpath counts; avoiding: subtract.
\cue{Without replacement.} \(K\) special among \(N\), draw \(n\):
\[P(X=k)=\frac{{K\choose k}{N-K\choose n-k}}{{N\choose n}}.\]
\(\E X=nK/N\); \(\V X=np(1-p)(N-n)/(N-1)\), \(p=K/N\).
\cue{Events.} \(P(A^c)=1-P(A)\);\par
\(P(A\cup B)=P(A)+P(B)-P(A\cap B)\).
At least one: complement of none. Exactly one of two: \(P(A)+P(B)-2P(AB)\).
\cue{Conditioning.} \(P(A\mid B)=P(AB)/P(B)\).
\(P(AB)=P(A)P(B\mid A)\). For a partition \(B_i\):
\(P(A)=\sum_iP(B_i)P(A\mid B_i)\).
\[
 P(B_j\mid A)=\frac{P(B_j)P(A\mid B_j)}{\sum_iP(B_i)P(A\mid B_i)}.\]
Independent: \(P(AB)=P(A)P(B)\); disjoint positive-probability events are dependent. Pairwise independence need not imply mutual independence.
\head[early]{2\quad One RV: normalize, integrate, transform}
PMF: \(\sum_x f(x)=1\). PDF: \(\int f(x)\,dx=1\). Always \(f\ge0\); a density can exceed 1.
\(F(x)=P(X\le x)\); continuous: \(F'=f\), \(P(X=a)=0\).
\(P(a<X\le b)=F(b)-F(a)\).
Integer \(X\): \(P(X\ge k)=1-F(k-1)\), \(P(X>k)=1-F(k)\). Discrete mass is a CDF jump.
Valid CDF: nondecreasing, right-continuous, limits \(0,1\). General endpoints: \(P(a\le X\le b)=F(b)-F(a^-)\).
\cue{Unknown constant.} Sketch support, normalize, integrate over support \(\cap\) event.
\cue{Expectation.} \(\E g(X)=\sum g(x)f(x)\), or \(\int g(x)f(x)dx\). \(\V X=\E X^2-(\E X)^2\).
\(\E(aX+b)=a\E X+b\), \(\V(aX+b)=a^2\V X\).
Indicator: \(\E I_A=P(A)\); counts \(N=\sum I_i\) give \(\E N=\sum P(A_i)\), even if dependent.
\cue{Transforms.} Map discrete values; add masses with the same image. Piecewise PDF: split integrals at formula changes.
\cue{Extremes, i.i.d.} \(F_M(x)=F(x)^n\), \(F_m(x)=1-[1-F(x)]^n\). Differentiate for PDFs.
\head[early]{3\quad Joint RVs: support matters}
\(f_X(x)=\sum_y f(x,y)\) or \(\int f(x,y)dy\).
\(f_{Y|X}(y|x)=f(x,y)/f_X(x)\), if \(f_X(x)>0\).
\(\E[g(X,Y)]=\sum_x\sum_y g(x,y)f(x,y)\), or double integral.
\cue{Triangular support.} \(0<x<y<1\): \(f_X=\int_x^1f\,dy\), \(f_Y=\int_0^yf\,dx\). Fixing a value also restricts conditional support.
\cue{Independence.} \(f(x,y)=f_X(x)f_Y(y)\) everywhere; positive support must be a product set.
\col
\head[early]{4\quad Joint moments \& distribution recognition}
\[\Cov(X,Y)=\E(XY)-\E X\E Y,\quad
 \rho=\frac{\Cov(X,Y)}{\sigma_X\sigma_Y}.\]
\[\V(aX+bY)=a^2\sigma_X^2+b^2\sigma_Y^2+2ab\Cov(X,Y).\]
\(\Cov(aX+b,cY+d)=ac\Cov(X,Y)\). Independence implies covariance zero; converse false. Variances of independent sums \emph{and differences} add.
\cue{Conditional moments.} \(\E(Y|X=x)=\sum_y y f_{Y|X}(y|x)\); integrate for continuous \(Y\).\par
\(\E Y=\E[\E(Y|X)]\); \(\V Y=\E[\V(Y|X)]+\V[\E(Y|X)]\).
\head[early]{5\quad Distribution cues \quad (\(q=1-p\))}
\cue{Bernoulli: one yes/no.} \(P(X=1)=p\); mean \(p\), variance \(pq\).
\cue{Binomial: fixed \(n\), i.i.d. trials.}
\(P(X=k)={n\choose k}p^kq^{n-k}\); mean \(np\), variance \(npq\). At least one: \(1-q^n\).
\cue{Geometric: trials to first success.}
\(P(X=k)=q^{k-1}p\), \(k\ge1\); \(P(X>k)=q^k\) for integer \(k\ge0\). Mean \(1/p\), variance \(q/p^2\).
\cue{Negative binomial: trials to \(r\)th success.}
\(P(X=k)={k-1\choose r-1}p^rq^{k-r}\), \(k\ge r\). Mean \(r/p\), variance \(rq/p^2\).
\trap{If counting failures, subtract \(r\) from trials; the mean changes, the variance does not.}
\cue{Poisson: arrivals/count in interval.}
\(P(X=k)=e^{-\lambda}\lambda^k/k!\); mean = variance = \(\lambda\). Rate \(r\), duration \(t\): \(\lambda=rt\). Disjoint increments independent; independent Poisson counts add.
\cue{Uniform \(U(a,b)\): constant density.}
\(f=1/(b-a)\); \(\E X=(a+b)/2\), \(\V X=(b-a)^2/12\). Probability = interval overlap length / \((b-a)\).
\cue{Exponential: waiting for next arrival.}
\(f(x)=\lambda e^{-\lambda x}\), \(x\ge0\); survival \(e^{-\lambda x}\); mean \(1/\lambda\), variance \(1/\lambda^2\).
\(P(X>s+t|X>s)=P(X>t)\) (also geometric for integer times). Independent exponential minimum has rate \(\sum\lambda_i\).
\cue{Normal \(N(\mu,\sigma^2)\).} \(Z=(X-\mu)/\sigma\).
\(P(a<X<b)=\Phi((b-\mu)/\sigma)-\Phi((a-\mu)/\sigma)\).
Quantile with upper tail \(\gamma\): \(x=\mu+\sigma z_\gamma\). Independent normal linear combinations are exactly normal.
\cue{Approximations.} Binomial \(\approx\operatorname{Pois}(np)\) for rare successes. Binomial \(\approx N(np,npq)\) when \(np,nq\) sufficiently large.
\cue{Continuity correction (integer counts).}
\(P(X\le k)\to P(Y<k+.5)\);\par
\(P(X\ge k)\to P(Y>k-.5)\);\par
\(P(a\le X\le b)\to P(a-.5<Y<b+.5)\).
\trap{Do not continuity-correct ordinary continuous measurements. Normal notation gives variance, not SD.}
\head[early]{6\quad Combined probability question patterns}
\cue{Poisson given total.} Independent \(X\sim\operatorname{Pois}(\lambda)\), \(Y\sim\operatorname{Pois}(\eta)\):
\(X\mid X+Y=m\sim\operatorname{Bin}(m,\lambda/(\lambda+\eta))\).
\cue{Probability first, count second.} Find \(p=P(\text{one item meets threshold})\) using its normal/exponential model; number among \(n\) independent items is \(\operatorname{Bin}(n,p)\).
\cue{At least one with probability \(>c\).}
\(1-(1-p)^n>c\Rightarrow n>\log(1-c)/\log(1-p)\). Strict inequality: smallest integer \emph{greater} than bound.
\col
\head{7\quad Sampling \& CLT\hfill SECOND HALF}
\cue{Random sample.} \(X_1,\ldots,X_n\) i.i.d.; joint density/PMF \(\prod_i f_X(x_i)\).
SRS without replacement: each subset has probability \(1/{N\choose n}\); selections are dependent.
\cue{Enumerate a sampling distribution.} List ordered samples, compute each mean, then combine equal means by adding their sample probabilities. Distinct means are not necessarily equally likely.
\[\Xb=\frac1n\sum_iX_i,\qquad
 S^2=\frac{\sum_i(X_i-\Xb)^2}{n-1}.\]
\[\E\Xb=\mu,\quad\V\Xb=\frac{\sigma^2}{n},\quad
 \SE(\Xb)=\frac{\sigma}{\sqrt n}.\]
\(\sum X_i\): mean \(n\mu\), variance \(n\sigma^2\), SD \(\sigma\sqrt n\).
\cue{Exact vs approximate.} Normal population \(\Rightarrow\Xb\sim N(\mu,\sigma^2/n)\) for \emph{every} \(n\). Otherwise finite-variance i.i.d. CLT gives
\[\frac{\Xb-\mu}{\sigma/\sqrt n}\approx N(0,1),\qquad
 \frac{\sum X_i-n\mu}{\sigma\sqrt n}\approx N(0,1).\]
Course rule: large \(n\ge30\); strong skew/heavy tails may need larger \(n\). CLT makes the \emph{mean} approximately normal, not the individual observations.
\cue{LLN.} \(\Xb\xrightarrow{P}\mu\): averages concentrate near \(\mu\). This alone does not assert a normal shape. Chebyshev: \(P(|\Xb-\mu|\ge e)\le\sigma^2/(ne^2)\).
\head{8\quad CLT question patterns}
\cue{Probability of an average.}
\[P(a<\Xb<b)\approx
 \Phi\!\left(\frac{b-\mu}{\sigma/\sqrt n}\right)
 -\Phi\!\left(\frac{a-\mu}{\sigma/\sqrt n}\right).\]
\textbf{Sum tail:} \(P(\sum X_i>b)\approx1-\Phi((b-n\mu)/(\sigma\sqrt n))\). \textbf{W9 pre1:} \(n=100,\mu=40,\sigma=\sqrt{100}=10,b=4150\). \(z=(4150-100\cdot40)/(10\sqrt{100})=1.5\); tail \(.0668\).
\cue{Choose \(n\) for precision.}
\[P(|\Xb-\mu|\le e)\gtrsim1-\alpha
 \ \Longrightarrow\ n\ge\left(\frac{z_{\alpha/2}\sigma}{e}\right)^2.\]
For \(P(\Xb>\mu+d)\lesssim\alpha\), \(d>0\):
\[n\ge(z_\alpha\sigma/d)^2.\]
Round up after solving; these are CLT approximations.
\cue{SE target only.} \(\SE\le e\Rightarrow n\ge(\sigma/e)^2\): no \(z\) factor. SE multiplied by \(r\) needs sample size multiplied by \(1/r^2\); quarter SE needs \(16n\).
\head{9\quad Calculator-friendly data summaries}
Store \(A=\sum x_i\), \(B=\sum x_i^2\); then
\[\xb=A/n,\qquad s^2=\frac{B-A^2/n}{n-1}.\]
\(s=\sqrt{s^2},\quad\widehat{\SE}=s/\sqrt n\).
Expansion: \(\sum(x_i-\xb)^2=\sum x_i^2-n\xb^2\).
\trap{\(s^2\) is variance; \(s\) is SD; \(s/\sqrt n\) is estimated SE. If given variance, do not square it again. Use calculator sample SD (\(s_x\), divisor \(n-1\)).}
\cue{MRQ traps (W8 post1,4).} For \(Z=\sqrt n(\Xb-\mu)/\sigma\), \(\V Z=1\) exactly; \(\V(2Z)=4\). Omitting centering changes the mean, not the variance. CLT: \(P(\Xb>\mu)\approx.5\), even as the mean concentrates around \(\mu\).
\col
\head{10\quad Sampling pivots: normal population}
For an i.i.d. normal sample, \(\Xb\) and \(S^2\) are independent and
\[Z=\frac{\Xb-\mu}{\sigma/\sqrt n}\sim N(0,1),\quad
 Q=\frac{(n-1)S^2}{\sigma^2}\sim\chi^2_{n-1}.\]
\[T=\frac{\Xb-\mu}{S/\sqrt n}\sim t_{n-1}.\]
Independent standard-normal squares sum to \(\chi^2\), df = number of squares. For independent \(Z\sim N(0,1)\), \(U\sim\chi^2_\nu\): \(Z/\sqrt{U/\nu}\sim t_\nu\). Student \(t\): symmetric, heavier tails than normal, approaches normal as df grows.
\cue{Two independent normal samples.}
\[\frac{S_1^2/\sigma_1^2}{S_2^2/\sigma_2^2}
 \sim F_{n_1-1,n_2-1}.\]
If population variances equal, ratio is \(S_1^2/S_2^2\). Numerator df first. Definition: \((U/r)/(V/s)\sim F_{r,s}\) for independent chi-squares \(U,V\) with df \(r,s\).
Raw ratio \(U/V\sim(r/s)F_{r,s}\). Divide each chi-square by \emph{its own df} first. \(\E F_{r,s}=s/(s-2)\), \(s>2\): mean depends on denominator df. Ratios sharing a random denominator need not be independent.
\[F\sim F_{r,s}\ \Rightarrow\ 1/F\sim F_{s,r}.\]
\[P(F_{r,s}>c)=P(F_{s,r}<1/c).\]
As both df grow, \(\chi^2_r/r\xrightarrow{P}1\) since its variance is \(2/r\); hence \(F_{r,s}\xrightarrow{P}1\) (T9 Q6).
\cue{Chi-square normal approximation (W9 post2).}
\(\nu=n-1,\ \E Q=\nu,\ \V Q=2\nu\); large \(\nu\): \((Q-\nu)/\sqrt{2\nu}\approx N(0,1)\). For \(S^2>k\sigma^2\), use \(z=(\nu k-\nu)/\sqrt{2\nu}\). Example \(n=51,k=1.4\): \(z=(50\cdot1.4-50)/\sqrt{100}=2\); tail \(1-\Phi(2)\approx.0228\), not exact.
\head{11\quad Quantiles \& pivot probability questions}
\textbf{All critical values here use upper-tail notation:}\par
\(P(Z>z_\gamma)=\gamma\), \(P(T_\nu>t_{\nu;\gamma})=\gamma\); similarly \(\chi^2_{\nu;\gamma}\), \(F_{r,s;\gamma}\).
\cue{Variance probability.} \(P(S^2>c)=P(Q>(n-1)c/\sigma^2)\). For \(P(S^2<c)=p\):
\[c=\frac{\sigma^2\chi^2_{n-1;1-p}}{n-1}.\]
\textbf{Example:} \(n=13,p=.2\), \(\chi^2_{12;.8}=7.8073\): \(c=.6506\sigma^2\).
\cue{Random \(S\) in the bound.}
\(P(|\Xb-\mu|<cS)=1-\alpha\Rightarrow c=t_{n-1;\alpha/2}/\sqrt n\).
\textbf{Example:} \(n=25\): \(c=2.064/5=.4128\).
\trap{\(t\) and chi-square pivots are exact because of normality, not merely because \(n\) is small. Reciprocal \(F\) swaps df \emph{and} reverses the inequality.}
\head{12\quad Point estimation: bias \& comparison}
Estimator = statistic (random); estimate = its observed value.
\(\operatorname{Bias}(T)=\E T-\theta\); unbiased iff \(\E T=\theta\).
\(\operatorname{MSE}(T)=\V T+[\operatorname{Bias}(T)]^2\).
Among unbiased estimators, smaller variance is preferable.
\cue{Correct a bias.} If \(\E T=a\theta+b\), use \((T-b)/a\) (known \(a\ne0,b\)).
\cue{Centering costs one df.} \(\sum(X_i-\Xb)^2=\sum(X_i-\mu)^2-n(\Xb-\mu)^2\); expectation \((n-1)\sigma^2\). Thus \(S^2\) unbiased even without normality; \(S\) generally biased. Under normality, division by \(\sigma^2\) gives \(\chi^2_{n-1}\) with estimated mean, \(\chi^2_n\) with true mean.
\end{multicols*}
\newpage
\begin{multicols*}{4}
\head{13\quad Choose the formula FIRST\hfill 2 / 2}
\textbf{Symbols:} \(\alpha=1-\text{confidence level}\), \(\nu=\text{df}\); \(E=\text{margin of error}\), \(c=\text{critical value}\), \(\SE=\text{standard error}\).
\textbf{1. Target:} one mean \(\mu\), difference \(\delta=\mu_1-\mu_2\), or paired mean \(\mu_D\)?\par
\textbf{2. Paired?} Same subjects/matched units \(\to D_i\), then use one-sample formulas.\par
\textbf{3. Variance known?} Given \(\sigma\) or \(\sigma^2\) \(\to z\). Given \(s\) or \(s^2\) \(\to\) unknown.\par
\textbf{4. Check assumptions:} normality, equal variances, and the specified \(z\) or \(t\) reference (cases below).\par
\textbf{5. Task:} CI uses estimate \(\pm\) critical \(\times\SE\); test uses (estimate\(-\)null)/\(\SE\).
\head{14\quad One mean: CI \& test toolkit}
\[\text{CI: }\xb\pm c\SE,\qquad
 \text{test statistic: }\frac{\xb-\mu_0}{\SE}.\]
\cue{Known population SD \(\sigma\).}
\(\SE=\sigma/\sqrt n\), \(c=z_{\alpha/2}\). Reference \(N(0,1)\): exact for normal population, approximate for large samples otherwise.
\cue{Unknown \(\sigma\), normal population (any \(n\)).}
\(\SE=s/\sqrt n\), \(c=t_{n-1;\alpha/2}\), reference \(t_{n-1}\).
\cue{Unknown \(\sigma\), explicit large-sample method.}
\(n\ge30\), \(\SE=s/\sqrt n\), \(c=z_{\alpha/2}\), approximate normal reference. For a normal population, \(t\) remains exact at any \(n\).
\trap{Quiz gives normal population + unknown variance + \(t\) cutoff? Use \(t\) at \(n=36\) or \(49\) too. \(s=\sigma\) numerically does not make \(\sigma\) known. Follow the stated reference method.}
\cue{Small and nonnormal?} These assumptions do not justify an exact \(t\) interval/test automatically.
\head{15\quad Two independent means}
\(n_i\): group \(i\)'s sample size; \(\xb_i\): sample mean; \(s_i\): sample SD; \(s_i^2\): sample variance. \(\sigma_i^2\): population variance. Target \(\delta=\mu_1-\mu_2\); estimate \(\hat\delta=\xb_1-\xb_2\); null value \(\delta_0\).
\[\text{CI: }\hat\delta\pm c\SE,\qquad
 \text{test statistic: }\frac{\hat\delta-\delta_0}{\SE}.\]
\cue{Known variances.}
\(\SE=\sqrt{\sigma_1^2/n_1+\sigma_2^2/n_2}\); use \(z\). Exact for normal populations; CLT if both samples large.
\cue{Unknown, unequal/unrestricted; both large.}
\(\SE=\sqrt{s_1^2/n_1+s_2^2/n_2}\); approximate \(z\).
\cue{Unknown, equal population variances: pool.}
\textbf{1. Pooled sample variance:}
\[s_p^2=\frac{(n_1-1)s_1^2+(n_2-1)s_2^2}{n_1+n_2-2}.\]
\textbf{2. SD then SE of the difference:}
\[s_p=\sqrt{s_p^2},\qquad
 \SE=s_p\sqrt{\frac1{n_1}+\frac1{n_2}}.\]
\textbf{3. CI or test:} \(\hat\delta\pm c\SE\), or \(t=(\hat\delta-\delta_0)/\SE\). If given SDs, square them in step 1; if given variances, use them directly.
Normal populations: \(t\), df \(n_1+n_2-2\) (exact, any sizes). Course uses \(z\) when both samples large. Use the equal-variance formula when explicitly requested.
\trap{For large-sample \(z\), both groups must be large; a large total alone is insufficient.}
\head{16\quad Critical values to recall}
\begin{tabular}{@{}lrrrr@{}}
Confidence &90\%&95\%&98\%&99\%\\
\(\alpha/2\)&.05&.025&.01&.005\\
\(z_{\alpha/2}\)&1.645&1.960&2.326&2.576
\end{tabular}\par
\(\Phi(z_\gamma)=1-\gamma\); \(\Phi(-z)=1-\Phi(z)\).
\col
\head{17\quad Paired observations: reduce to one sample}
\cue{Recognize.} Same patient/task before and after; matched pairs; same units under two conditions. Independence is needed \emph{between pairs}, not within each pair.
Define direction explicitly: e.g. \(D=\text{before}-\text{after}\) for time reduction. Then reduction means \(\mu_D>0\).
\(d_i\): observed difference for pair \(i\); \(\db\): mean difference; \(s_D\): SD of differences; \(\mu_{D0}\): null mean difference; \(n\): number of pairs.
\[\db=\frac{\sum d_i}{n},\quad
 s_D^2=\frac{\sum d_i^2-(\sum d_i)^2/n}{n-1}.\]
\[\text{CI: }\db\pm c\frac{s_D}{\sqrt n},\qquad
 T=\frac{\db-\mu_{D0}}{s_D/\sqrt n}.\]
Normal \emph{differences}: exact \(t_{n-1}\); explicit large-sample method: approximate \(z\). \(n\) counts pairs. Normal marginals alone do not guarantee normal differences for dependent pairs. Similar observed variances do not establish equal population variances for pooling.
\cue{Only separate summaries and correlation?}
Here \(r\) is the sample correlation of the paired \(X,Y\) values.
\[\db=\xb-\bar y,\qquad s_D^2=s_X^2+s_Y^2-2rs_Xs_Y.\]
\trap{Do not pool paired data. Do not use \(s_X^2+s_Y^2\) unless the within-pair covariance is zero. Reversing \(D\) flips the mean, statistic, and alternative direction, but not the SD.}
\head{18\quad Confidence interval question patterns}
\cue{Build an interval.} Identify case \(\to\) compute \(\SE\) \(\to\) choose \(c\) \(\to E=c\SE\) \(\to\) estimate \(\pm E\).
\cue{Work backwards from \((L,U)\).}
\[\hat\theta=\frac{L+U}{2},\quad E=\frac{U-L}{2},\quad
 \SE=\frac{E}{c}.\]
Then test any null using \(w=(\hat\theta-\theta_0)c/E\). Reverse target sign: interval for \(-\theta\) is \((-U,-L)\).
One sample or paired: \(s=\SE\sqrt n\) if estimated SE is used. E.g. 95\% \((.58,4.50)\): center \(2.54\), \(E=1.96\), \(\SE=1\).
\cue{Required \(n\), known/planning SD.}
\[n_{\min}=\left\lceil\left(\frac{z_{\alpha/2}\sigma}{E}\right)^2\right\rceil.\]
If total length is \(W\), first set \(E=W/2\). A pilot SD gives a planning approximation, not a guaranteed known-variance bound.
\textbf{Example:} \(\sigma=8,W=2.5\), 95\%: \(E=1.25\), \(n=\lceil157.351936\rceil=158\).
\cue{Change sample size/confidence (fixed SD).}
\[\frac{E_{\rm new}}{E_{\rm old}}=
 \frac{c_{\rm new}}{c_{\rm old}}\sqrt{\frac{n_{\rm old}}{n_{\rm new}}}.\]
Fixed confidence: width reduced by \(r\) means \(n_{\rm new}=\lceil n_{\rm old}/(1-r)^2\rceil\). A 25\% reduction: \(81/.75^2=144\). Half width needs four times \(n\). Formula assumes critical value fixed (e.g. \(z\)).
\cue{Confidence change, same data.}
\(\%\Delta E=100(c_{\rm new}/c_{\rm old}-1)\). 95\%\(\to\)99\% normal CI: \(31.43\%\) wider.
If variance is multiplied by \(a\) and \(n\) by \(b\), width multiplier is \((c_{\rm new}/c_{\rm old})\sqrt{a/b}\). Triple both variance and \(n\): unchanged width. With one sample size fixed and the other \(\to\infty\), two-mean SE tends to the fixed group's SE, not zero.
\cue{Interpret.} A 95\% CI procedure covers the fixed true parameter in 95\% of repeated samples under its assumptions. It does not contain 95\% of observations. Higher confidence \(\to\) wider interval; larger \(n\) \(\to\) narrower interval.
\col
\head{19\quad Hypothesis tests: choose the TAIL}
Let \(\theta\) be \(\mu,\delta\), or \(\mu_D\). Null value is \(\theta_0\); statistic \(w=(\hat\theta-\theta_0)/\SE\). Select \(z\) or \(t_\nu\) using sections 13--17.
\cue{Right: ``greater / exceeds / increased''.}
\(H_1:\theta>\theta_0\). Reject if \(w>c_\alpha\).\par
\(p=P(W\ge w)=1-G(w)\).
\cue{Left: ``less / below / reduced''.}
\(H_1:\theta<\theta_0\). Reject if \(w<-c_\alpha\).\par
\(p=P(W\le w)=G(w)\).
\cue{Two-sided: ``differs / changed / unequal''.}
\(H_1:\theta\ne\theta_0\). Reject if \(|w|>c_{\alpha/2}\).\par
\(p=2[1-G(|w|)]=2\min\{G(w),1-G(w)\}\).
Here \(G=\Phi\) for \(z\), or the \(t_\nu\) CDF; \(c_\gamma=z_\gamma\) or \(t_{\nu;\gamma}\).
\trap{One-sided \(p\) uses the signed statistic. Absolute value is only for the symmetric two-sided rule. Use the direction of \(H_1\), not the sign you hoped to obtain.}
\head{20\quad Reject only in the rejection region}
\textbf{Two-sided:}\par
\centerline{\textbf{reject}\quad\(\big|\,-c\quad0\quad c\,\big|\)\quad\textbf{reject}}
\centerline{\(w<-c\)\hfill\(-c\le w\le c\)\hfill\(w>c\)}
Between \(-c\) and \(+c\): \textbf{do not reject}.\par
\textbf{Anchor:} \(-2.306<1.76<2.306\), so do not reject. A positive statistic alone proves nothing.
\cue{p-value decision.} Reject for \(p<\alpha\) here; if the question specifies \(\le\), use it. Match critical-value inequalities to that boundary rule.
\cue{What p means.} Under \(H_0\), probability of a statistic at least as extreme in the direction(s) of \(H_1\). It is not \(P(H_0\mid\text{data})\).
\cue{Composite null.} For \(H_0:\theta\le\theta_0\) against \(>\), compute at boundary \(\theta_0\); similarly for \(\ge\) against \(<\).
\cue{Smallest \(n\) to reject: known \(\sigma\).} For a fixed observed \(\xb>\mu_0\), right-tail test:
\[\frac{(\xb-\mu_0)\sqrt n}{\sigma}>z_\alpha
\ \Rightarrow\ n>\left(\frac{z_\alpha\sigma}{\xb-\mu_0}\right)^2.\]
For two-sided tests use \(z_{\alpha/2}\) and \(|\xb-\mu_0|\). Minimum integer: \(\lfloor\text{bound}\rfloor+1\) for strict \(>\). W11 post9: \((2.576\cdot4/.8)^2=165.8944\Rightarrow166\). This fixes the observed mean; it is not power.
\cue{Conclude in context.} Reject: sufficient evidence for \(H_1\) at level \(\alpha\). Do not reject: insufficient evidence for \(H_1\); this does not prove \(H_0\).
\head{21\quad CI--test duality \& endpoint tricks}
A matching two-sided \(100(1-\alpha)\%\) CI excludes \(\theta_0\) iff the two-sided level-\(\alpha\) test rejects (same method, strict convention).
\cue{At an exact endpoint.} Two-sided \(p=\alpha\). At the upper endpoint, left-tail \(p=\alpha/2\), right-tail \(p=1-\alpha/2\); swap at lower endpoint. Rounded endpoints give approximate equality.
\textbf{96\% CI \([18,24]\), test \(\mu_0=24\):}\par
\(\xb=21\); \(p_<=.02,\ p_>=.98,\ p_{\ne}=.04\). Under strict \(p<.04\), only the left test rejects.
\cue{One-sided vs two-sided.} A 95\% two-sided CI can contain zero while a one-sided 5\% test rejects: its cutoff is smaller. A matching 95\% \emph{one-sided bound}, or 90\% two-sided interval in the relevant direction, matches that one-sided test.
\cue{Bound \(p\) using nested CIs (W12 post8).} Null outside 90\% but inside 95\% CI \(\Rightarrow .05<p<.10\); do not reject at 5\%. Compare with the actual null value, not automatically zero.
\col
\head{22\quad Quiz patterns: pivots, pooling, correlation}
\cue{Pooled test, nonzero null (W12 pre9).}
\(n_1=8,n_2=14,\ s_1=8,s_2=4\); observed difference \(57.3-50=7.3\). Test \(H_1:\delta>3\).
\textbf{Formula first (same steps as section 15):}
\[s_p^2=\frac{(n_1-1)s_1^2+(n_2-1)s_2^2}{n_1+n_2-2},\]
\[\SE=\sqrt{s_p^2\!\left(\frac1{n_1}+\frac1{n_2}\right)},\qquad
 t=\frac{(\xb_1-\xb_2)-\delta_0}{\SE}.\]
\textbf{Now substitute:} \(s_p^2=[7(8^2)+13(4^2)]/20=32.8\); \(\SE=\sqrt{32.8(1/8+1/14)}=2.5383\). Then \(t=(7.3-3)/2.5383=1.694\), df \(8+14-2=20\).\par
At 5\% right tail: \(1.694<1.725\), do not reject.
\cue{Paired correlation threshold (W12 pre10).}
\textbf{General:} let \(a=\db-\mu_{D0}>0\) and \(q=t_{n-1;\alpha}>0\) for a right-tail test. Then
\[t=\frac{a\sqrt n}{\sqrt{s_X^2+s_Y^2-2rs_Xs_Y}}>q\]
\[\Longleftrightarrow\quad r>\frac{s_X^2+s_Y^2-na^2/q^2}{2s_Xs_Y}.\]
\textbf{Substitute:}
\(n=16,\db=4.5,\mu_{D0}=2,s_X=8,s_Y=6\).
\(a=4.5-2=2.5,\ q=1.753\); \(t=10/\sqrt{100-96r}\). Thus \(r>[100-16(2.5)^2/1.753^2]/96=.7027\).
Higher positive \(r\) reduces difference variance. Check \(-1\le r\le1\), positive denominator and numerator sign before squaring; division by a negative flips an inequality.
\cue{Squared \(t\) (W9 post1).} \(W=n(\Xb-\mu)^2/S^2=T_{n-1}^2\sim F_{1,n-1}\). Upper-tail \(\alpha\) cutoff: \(t_{n-1;\alpha/2}^2\). \(n=16\): \(2.131^2=4.5412\), not \(2(2.131)\).
\cue{Weighted estimator (W9).} For i.i.d. data, \(T=\sum a_iX_i\): \(\E T=\mu\sum a_i\), \(\V T=\sigma^2\sum a_i^2\). Unbiased for every \(\mu\) iff \(\sum a_i=1\). Subset mean of \(m\) values: unbiased, variance \(\sigma^2/m\).
\textbf{Bias direction:} \(T=\sum X_i/(n+2)\) has bias \(-2\mu/(n+2)\), downward when \(\mu>0\).
\head{23\quad Critical cutoffs, Type II error \& power}
\(\mu_0\): null mean; \(\mu_a\): specified true mean; \(u,\ell\): upper/lower cutoffs in the original units.
\(\SE=\sigma/\sqrt n\). Right-tail rejection cutoff: \(u=\mu_0+z_\alpha\SE\); left-tail: \(\ell=\mu_0-z_\alpha\SE\). Two-sided: use \(z_{\alpha/2}\) in both cutoffs.
\cue{At true mean \(\mu_a\), keep the cutoff fixed.}
\[\beta_{\rm right}=\Phi\!\left(\frac{u-\mu_a}{\SE}\right),\quad
 \beta_{\rm left}=1-\Phi\!\left(\frac{\ell-\mu_a}{\SE}\right).\]
Two-sided: \(\beta=\Phi((u-\mu_a)/\SE)-\Phi((\ell-\mu_a)/\SE)\). Power \(=1-\beta\). Exact for normal data with known \(\sigma\); CLT approximation otherwise.
\textbf{W11 pre10, substitute:} \(\mu_0=100,\sigma=12,n=36,\alpha=.05\): \(\SE=12/\sqrt{36}=2\); \(u=100+1.645(2)=103.29\). At \(\mu_a=104\): \(\beta=\Phi((103.29-104)/2)=.3613\), power \(1-.3613=.6387\).
\cue{Error traps.} Type I = reject a true null (level \(\alpha\)); Type II = retain a false null (\(\beta\), depends on \(\mu_a\)). Power is not \(1-\alpha\). Choosing the one-sided direction \emph{after} seeing the sign gives actual Type I error \(2\alpha\) for this symmetric rule (W11 pre9).
\textbf{Wrong-direction tail (W12 post7):} \(H_1:\mu<50\), but \(z=+1.75\): \(p=\Phi(1.75)=.9599\), not \(.0401\). Do not reject at 5\%.
\textbf{Before committing:} target/sign \(\to\) paired? \(\to\sigma\) vs \(s\) \(\to\SE\) \(\to\) df \(\to\) tail \(\to\) inequality \(\to\) contextual conclusion. Round sample size up and final answers last.
\end{multicols*}
\end{document}
